\documentclass[runningheads,a4paper]{llncs}
\usepackage[shortend,ruled]{algorithm2e}
\usepackage{amsmath,amsfonts}

\title{Column Generation in Boolean Optimization}
\author{Jordi Planes}
\institute{Universitat de Lleida}

\bibliographystyle{splncs03}

\begin{document}

\maketitle

\begin{abstract}
Recently several MaxSat algorithms based on SAT reduction have
appeared in the literature. Our paper attemps to present a general
methodology , gathering all the algorithms in a common framework.
%A common framework to demonstrate the soundness of the MaxSat unsat-based solvers.
%We show that several recent algorithms applied to the MaxSat solving are the Benders decomposition~\cite{benders62} when applying the inference dual~\cite{hooker99}. We discuss the several algorithms and how are related to the Benders decomposition. This leads to show several features (incrementality) of the algorithms and other possible algorithms.

%El fet que el dual de MaxSat sigui el dual de Hooker~\cite{hooker99} fa que siguin més
%fàcils d'explicar els algorismes SAT-oracle based. Son l'aplicacio del
%metode de Benders a la logica propositional~\cite{benders62} i
%canviant el metode Simplex pel metode DP~\cite{DP60}.
%Aixo ens permet entendre millor els metodes, verificar-ne la
%correctesa i millorar-los i trobar-ne de nous.

%Tambe demostrarem la no incrementalitat dels algorismes, tant SAT-UNSAT com UNSAT-SAT.
\end{abstract}

% Esquema bàsic (4 pages at most):
% \begin{itemize}
% \item dualitat s -- r
% \item Identificar inference dual (i/o altres)
% \item Metodes de generacio de columnes amb Const(k) i relacio unsat-sat
% \item Metodes sat-unsat
% \end{itemize}

\section{Introduction}\label{sec:intro}

TO BE DONE: Max-2-SAT dual = graph duality. MaxSat dual is MinOnes?
One can get it applying resolution.
Maxflow, Mincut theorem?

In operations research, there are a well known set of column
generation methods (e.g., Dantzig-Wolfe~\cite{dantzig-wolfe},
Benders cuts~\cite{benders62}) which are based on
the relaxation of an optimization problem into an easier problem
(usually a linear problem solved by the Simplex method) and
iteratively solve the initial problem.

Recently, a set of algorithms have been created to solve the MaxSat
problem. Such algorithms are original in the sense that they
iteratively reduce the MaxSat problem into Sat.

The complexity of MaxSat is $P^{NP[log n]}$ , meaning that it can be
solved with a logarithmic number of calls to a SAT oracle~\cite{papadimitriou94}.

In this paper we extend the idea of inference dual~\cite{hooker99} to
the MaxSat problem.


\section{Preliminaries}

Given a CNF instance $\phi$, the MaxSat problem is the problem of finding an
assignment which maximizes the number of satisfied clauses in $\phi$.

The MaxSat problem can be written as
\begin{eqnarray}
\min & \sum s_i \\
\text{subject to} & s_i + \sum a_{ij} x_i \ge 1 \\
& a_{ij}, x_i, s_i \in \{ 0, 1 \}
\end{eqnarray}
where $x_i$ is a Boolean variable in the CNF and $a_{ij}$ is the
selector variable for variable $x_i$ in clause $j$ and $s_j$ is the
selector of clause $j$.
Minimize the number of selectors in order to make the formula satisfiable.

The dual of the problem above is:
\begin{eqnarray}
\min &\sum \sigma_i \\
\text{subject to} & \sigma_i + \prod \alpha_{ij} c_i = 1 \\
& \alpha_{ij}, \sigma_i, c_i \in \{ 0, 1 \}  
\end{eqnarray}


\begin{eqnarray}
\min &\sum \sigma_i \\
\text{subject to} & \sigma_i + F( \alpha_{ij} c_i ) \le 1 \\
& \alpha_{ij}, \sigma_i, c_i \in \{ 0, 1 \}  
\end{eqnarray}

\begin{eqnarray}
\max & \sum \sigma_i \\
\text{subject to} & \sigma_i + F( \alpha_{ij} c_i ) \le 1 \\
& \alpha_{ij}, \sigma, c_i, F( \alpha_{ij} c_i ) \in \{ 0, 1 \} 
\end{eqnarray}

where $\sigma_i$ is the selector of the proof $i$ and $\alpha_{ij}$ is the
selector of clause $i$ in proof $j$.
Minimize the number of proofs with condition $\le 1$ in order to make
the formula satisfiable.
$F( \phi  )$ is a function which returns 0 if the formula $\phi$ is
satisfiable, otherwise returns 1.

This problem corresponds to 
\begin{eqnarray}
\max & \sum s_i \\
\text{subject to} & \prod s_i c_i = 0
\end{eqnarray}
The minimal removal ($s_i = 0$) of $c_i$ which makes the formula satisfiable.

An optimization problem can be written 
\begin{equation}
\min \{ f(x) | \mathcal{C} \}
\label{eq:opt-prob}
\end{equation}
where $f(x)$ is a real-valued function, $\mathcal{C}$ is a constraint set containing variables $x=(x_1, \ldots, x_n)$, and $D$ is a domain of $x$. A solution $\bar x\in D$ is feasible when it satisfies $\mathbb{C}$ and is optimal when $f(\bar x) \le f(x)$ for all feasible $x$.

Given a CNF formula, the MaxSat is the problem of finding an assignment with the maximum number of satisfied clauses.
Using the notation in~\ref{eq:opt-prob}, the MaxSat problem can be written as
\begin{equation}
\min_{a,x,w\in \{ 0, 1 \}} \{ \sum w_j | \sum a_{ij} x_i + w_j  \ge 1 \}
\label{eq:maxsat}
\end{equation}

The inference dual of \ref{eq:opt-prob} is the problem of finding the greatest lower bound on $f(x)$ that can be inferred from $\mathcal{C}$ within a given proof system. The inference dual can be written
\begin{equation}
\max_{P\in \mathcal{P}} \{ v\ |\ \mathcal{C} \overset{P}{\vdash} (f(x) \ge v) \}
\label{eq:dual}
\end{equation}
where $\mathcal{C} \overset{P}{\vdash} (f(x) \ge v)$ indicates that proof $P$ deduces $f(x) \ge v$ from $\mathcal{C}$. The domain of variable $P$ is a family $\mathcal{P}$ of proofs. A pair $(\bar v, \bar P)$ is a feasible solution of \ref{eq:dual} if $P\in\mathcal{P}$ and $\mathcal{C} \overset{\bar P}{\vdash} (f(x) \ge \bar v)$, and it is optimal if $\bar v \ge v$ for all feasible $(v,P)$.

Equivalently, it can be defined as:
\begin{equation}
\min_{P\in \mathcal{P}} \{ v\ |\ \mathcal{C} \overset{P}{\vdash} (f(x) \le v) \}
\label{eq:dual2}
\end{equation}

%%%

\begin{lemma}
Given a MaxSat primal assignment, there is a unique MaxSat dual assignment.
\end{lemma}
\begin{proof}
Hitting set, unuseful variables, variables always in MUSes, \dots \cite{KLM09}
\end{proof}

%%%

\subsection{Models for Column Generation}

Barnhart et al.~\cite{bjnsv98} considere a general problem $P$ of the
form:

\begin{displaymath}
\begin{array}{rl}
\max & c(x) \\
& Ax \le b, \\
& x \in S,\\
& x \quad integer
\end{array}
\end{displaymath}

The function $c(x)$ is needed to be easy to solve with $Ax \le b$
removed, and $S^* \subseteq S $.
Its column generation form is:

\begin{displaymath}
\begin{array}{rl}
\max & \sum c_k \lambda_k \\
& \sum a_k \lambda_k \le b \\
& \sum \lambda_k = 1 \\
& \lambda_k \in \{ 0, 1 \} \quad k=1,\ldots,p
\end{array}
\end{displaymath}

\section{Related work}

Column generation in Integer Programming~\cite{bjnsv98,vw96}

Column generation in MaxSat~\cite{hjp98}

Column generation in Constraint Programming~\cite{FJKKSV02}

%%%

\section{Column generation in MaxSat}

Application of the Dantzig-Wolfe
decomposition to a structured optimization problem in which
formulations with an exponential number of variables are
obtained. Generating these variables on demand by means of delayed
column generation is identical to performing a cutting plane on the
respective dual problem.

%%%

\begin{definition}[MUS]
A MUS $\phi$ is a CNF minimally unsatisfiable.
\end{definition}

\begin{definition}[Relaxed CNF]
A relaxed CNF $\phi^r$ is a CNF $\phi$ with an additional (relaxing)
variable for every clause, i.e.:
$$\phi^r = \{ c_i \vee b_i : \forall c_i \in \phi \}.$$
\end{definition}

Notice that a relaxed CNF without any constraint to the variables
$b_i$ is satisfiable. 

\begin{definition}[Constrained Relaxed CNF]
A constrained relaxed CNF $\phi^r(k)$ is a relaxed CNF $\phi^r$ with
an additional constraint over the relaxing variables $\sum r_i = k$
where $k\in \mathbb{N}$.
\end{definition}

\begin{lemma}
Let $\phi^r(k)$ be a constrained relaxed CNF, $\phi$ be a CNF,
and $f(\phi)$ be the number of falsified
clauses in a MaxSat solution of $\phi$. Then,
 $$f(\phi^r(k)) = \max(  f(\phi) - k, 0 ). $$
\end{lemma}

\begin{proof}
A MaxSat solution of $\phi^r \cup (\sum b_ i= k)$
has $k$ different variables $b_i = 1$. Every of such
variables correspond to a falsified clause in a MaxSat
solution of $\phi$, since this will bring the largest number of
satisfied clauses.
If $k \ge f(\phi)$, all the falsified clauses are assigned with $b_i =
1$, so  $\phi^r \cup (\sum b_ i= k)$ is satisfiable.
\end{proof}

Notice that there is no need to relax the whole formula $\phi$. Any falsified clause in a MaxSat
solution belongs to a MUS. 
%So, only the set of clauses belonging to a MUS need to be relaxed.
%For any of the lemmata, notice that there is no need to relax the whole formula $\phi$ in
%any of the lemmata to be correct. 
The set of clauses to be relaxed can be reduced to the set of clauses
of $k$ MUSes.

Given $\phi^r_ 1 \cup (\sum b_ i= k)$, $\phi^r_1$ relaxing the whole
formula, and $\phi^r_ 2 \cup (\sum b_ i= k)$, $\phi^r_2$ relaxing $k$
MUSes which contain a clause with $b_i=1$, then both are MaxSat equivalent.

\begin{corollary}
Given $\phi$ a CNF, $k \in \mathbb{N}$, and $\phi^r$  a relaxed CNF with $\sum b_i = k$,
let $k_m$ be the minimum $k$ that makes the formula satisfiable, and
let $f(\phi)$ be the number of falsified clauses in a MaxSat
solution of $\phi$. Then, $k_m = f(\phi)$.
\label{lem:rcnf}
\end{corollary}

\begin{proof}
For any $k<f(\phi)$, there exists an falsified
clause in any MaxSat solution for $\phi^r$ with $b_i=0$, which makes $\phi^r$ unsatisfiable.
For any $k \ge f(\phi)$, there exists at least one combination of
falsified clauses
in any MaxSat sotlution with $b_i=1$ which makes the formula $\phi^r$ satisfiable.
\end{proof}

Notice that, given a MUS, $k_m = 1$.

\begin{lemma}[Composition of $\phi^r(k)$ is (sub)additive]
Given $\phi$ a CNF, $\phi^r_1$ a relaxed CNF $\phi$  with $\sum b_i =
k_1 + k_2$, and $\phi^r_2$ a relaxed CNF $\phi$  with $\sum b_i^1 =
k_1 \cup \sum b_i^2 = k_2$.
Then $\phi^r_1$ and $\phi^r_2$ are MaxSat equivalent.
\end{lemma}

\begin{proof}
Let $f(\phi)$ the number of falsified clauses in a MaxSat solution of
$\phi$.
Assuming $k_1 + k_2 < f(\phi)$, then $f(\phi^r_1) = f(\phi) -(k_1+k_2)$
and $f(\phi^r \cup \sum b_i^1 = k_1)= f(\phi) - k_1$.  
Adding constraint $\sum b_i^2 = k_2$ to formula  $\phi^r \cup \sum b_i^1 = k_1$
reduces the number of unsatisfied clauses by $k_2$, so 
$f(\phi^r \cup \sum b_i^1 = k_1 \cup \sum b_i^2 = k_2 ) = f(\phi) -  k_1 -
k_2$.
\end{proof}

\begin{corollary}
Given $\phi$ a CNF, and $\phi^{r(k)}$ a relaxed CNF with $\sum b_i^1 =
1 \cup \cdots \cup \sum b_i^k = 1$. Let $k_m$ be the minimum $k$ that makes the formula
$\phi^{r(k)}$ satisfiable, and let $f(\phi)$ be the number of falsified clauses in a
MaxSat solution of $\phi$. Then, $k_m = f(\phi)$.
\label{lem:rcnf2}
\end{corollary}

\begin{proof}
%At least, there are $f(\phi)$ MUSes in $\phi$. 
For any $k<f(\phi)$, there exists a set of falsified clauses in a MaxSat solution of $\phi^r$ with
$b_i=0$, which makes $\phi^r$ unsatisfiable.
For any $k \ge f(\phi)$, there exists at least one combination of
falsified clauses in a MaxSat solution of $\phi$
with $b_i=1$, which makes the formula $\phi^{r(k)}$ satisfiable.
\end{proof}

\section{Integer Programming}

$$
\begin{array}{ll}
\max \sum x_{ij}\\
& \sum a_{ij} + x_{ij} \le 1 \\
& \sum x_{ij} = 1 \\
a_{ij}, x_{ij} \in \{0,1\} 
\end{array}
$$

Relaxation
$$
\begin{array}{ll}
& \sum a_{ij} + x_{ij} \le 1 \\
& \sum x_{ij} = 1 \\
a_{ij}, x_{ij} \in \{0,1\} 
\end{array}
$$

Column generation
$$
\begin{array}{ll}
& \sum a_{ij} + x_{ij} + y_{ij}\le 1 \\
& \sum x_{ij} = 1 \\
& \sum y_{ij} = 1 \\
a_{ij}, x_{ij} \in \{0,1\} 
\end{array}
$$

\section{Unsatisfiablity based MaxSat algorithms}

Faltaria fer:
\begin{itemize}
\item IMPORTANT: Explicar algorismes MSU1, MSU3, \textbf{MSU4} i
  relacionar-los
\item IMPORTANT: Explicar algorismes anteriors a weighted i
  relacionar-los
\item Aplicar resolution (row generation) després de column generation
\item Identificar els casos on es pot aplicar local search unsat~\cite{AS07} en
  comptes de systematic: TOTS els casos.
\item Com es pot millorar Sat-Unsat: identificant clausules no necessaries (molt dificil)
\item Com extendre-ho a weighted, quines tecniques es podrien utilitzar
\end{itemize}

\bibliography{jordi,others}

\end{document}
